% 本文件由 LaTeX 排版 Agent 根据有效 translation.json 自动生成。
% author=Gregory Thaumaturgus; work_id=0611; title=论玛 6:22-23

\chapter{论\BibleRef{玛 6:22-23}}

% structure: p0001 source=h1 target=omitted reason=page-title-already-rendered-by-parent

\BibleQuote{身体的光就是眼睛；所以，如果你的眼是单纯的，你的全身就充满光明。但如果你眼是邪恶的，你的全身就充满黑暗。所以，如果你内在的光成为黑暗，那黑暗该是多么大啊！}\par

\OriginalTerm{单纯的眼}{single eye}就是\OriginalTerm{无伪的爱}{love unfeigned}；因为当身体被它光照时，它只能透过外在肢体表达与内在思想完全一致的事物。但\OriginalTerm{邪恶的眼}{evil eye}是假装的爱，也称为\OriginalTerm{伪善}{hypocrisy}，它使整个人身成为黑暗。我们必须考虑，仅适合\OriginalTerm{黑暗的工作}{work meet for darkness}可能存在于人内，而透过外在肢体，他可能产生看似属于光明的话语；因为确实有些人实质上是狼，尽管他们可能披着羊皮。这样的人只洗净杯盘的外面，却不明白，除非这些东西里面被洗净，外面本身也不能成为洁净。因此，为明确驳斥这样的人，救主说：\BibleQuote{如果你内在的光成为黑暗，那黑暗该是多么大啊！}这就是说，倘若那在你看来是光明的爱，因隐藏在你内的某些伪善，实际上是适合黑暗的工作，那么你明显的过犯又会怎样呢！\par

\SourceRecord{Translated by S.D.F. Salmond.；Ante-Nicene Fathers；Vol. 6.；Edited by Alexander Roberts, James Donaldson, and A. Cleveland Coxe.；Buffalo, NY: Christian Literature Publishing Co.,；1886.；<http://www.newadvent.org/fathers/0611.htm>.}
